\documentclass[10pt,a4paper]{amsart}
\usepackage{amsmath,amssymb,amsthm,graphicx}
\newcommand{\Sy}{\mathcal S}
\newcommand{\hh}{\mathbb{H}}
\newcommand{\D}{\mathcal{D}}
\newcommand{\J}{\mathcal{J}}
\newcommand{\Ab}{\mathcal{A}b}
\newcommand{\bbn}{\mathbb{N}}
\newcommand{\bbq}{\mathbb{Q}}
\newcommand{\R}{\mathbb{R}}
\newcommand{\RP}{\mathbb{RP}}
\newcommand{\CP}{\mathbb{CP}}
\newcommand{\HP}{\mathbb{HP}}
\newcommand{\GL}{\operatorname{GL}}
\newcommand{\Z}{\mathbb{Z}}
\newcommand{\ZZ}{\mathbb{Z}}
\newcommand{\Q}{\mathbb{Q}}
\newcommand{\QQ}{\mathbb{Q}}
\newcommand{\JJ}{\mathcal J}
\newcommand{\C}{\mathbb{C}}
\newcommand{\OO}{\mathbb{O}}
\newcommand{\HH}{\mathbb{H}}
\newcommand{\U}{\mathcal{U}}
\newcommand{\V}{\mathcal{V}}
\newcommand{\W}{\mathcal{W}}
\newcommand{\F}{\mathbb{F}}
\newcommand{\mbn}{\mathbf{n}}
\newcommand{\mbp}{\mathbf{p}}
\newcommand{\mbq}{\mathbf{q}}
\newcommand{\p}{\mathfrak p}
\newcommand{\q}{\mathfrak q}
\newcommand{\im}{\operatorname{Im}}
\newcommand{\mdc}{\operatorname{mdc}}
\DeclareMathOperator{\tr}{tr}
\DeclareMathOperator{\sgn}{sgn}
\DeclareMathOperator{\inter}{int}
\DeclareMathOperator{\Diff}{Diff}
\DeclareMathOperator{\ClImm}{ClImm}
\DeclareMathOperator{\Fun}{Fun}
\DeclareMathOperator{\coker}{coker}
\DeclareMathOperator{\PSh}{Psh}
\DeclareMathOperator{\Sp}{Sp}
\DeclareMathOperator{\Mod}{Mod}
\DeclareMathOperator{\Et}{\acute{E}t}
\DeclareMathOperator{\colim}{colim}
\DeclareMathOperator{\Hom}{Hom}
\DeclareMathOperator{\Ext}{Ext}
\DeclareMathOperator{\Ch}{Ch}
\DeclareMathOperator{\YExt}{YExt}
\DeclareMathOperator{\ev}{ev}
\DeclareMathOperator{\Supp}{Supp}
\DeclareMathOperator{\Tot}{Tot}
\DeclareMathOperator{\Spec}{Spec}
\DeclareMathOperator{\Tor}{Tor}
\DeclareMathOperator{\Der}{Der}
\DeclareMathOperator{\id}{id}
\DeclareMathOperator{\Gr}{Gr}
\DeclareMathOperator{\Op}{Op}
\DeclareMathOperator{\Emb}{Emb}
\DeclareMathOperator{\SympEmb}{SympEmb}
\DeclareMathOperator{\SympImm}{SympImm}
\DeclareMathOperator{\Th}{Th}
\DeclareMathOperator{\Map}{Map}
\DeclareMathOperator{\Aut}{Aut}
\DeclareMathOperator{\Id}{Id}
\DeclareMathOperator{\re}{Re}
\DeclareMathOperator{\Sq}{Sq}
\DeclareMathOperator{\End}{End}
\DeclareMathOperator{\Spin}{Spin}
\newcommand{\join}{\ast_{\Z/p}}
\newcommand{\cl}{\mathcal{C}{\rm l}}
\newcommand{\Cl}{\mathcal{C}{\rm l}}
\newcommand{\CCl}{\mathbb{C}{\rm l}}

\numberwithin{equation}{section}

\newtheorem{Theorem}{Theorem}[section]
\newtheorem*{Theorem*}{Theorem}
\newtheorem{Corollary}[Theorem]{Corollary}
\newtheorem{Lemma}[Theorem]{Lemma}
\newtheorem{Proposition}[Theorem]{Proposition}
 { \theoremstyle{definition}
\newtheorem{Definition}[Theorem]{Definition}
\newtheorem{Note}[Theorem]{Note}
\newtheorem{Example}[Theorem]{Example}
\newtheorem{Remark}[Theorem]{Remark} }
\title{Intersection number of two orthogonal almost complex structures on a closed Riemannian six-manifold in terms of their Chern classes}
\author{Gustavo Granja, Aleksandar Milivojević}
\date{doi:10.3842/SIGMA.2022.093; CC BY 4.0}
\begin{document}
\maketitle
\setcounter{section}{5}
\setcounter{Theorem}{6}
\noindent\textit{Editorial note.} All delivered mathematical text is the authors' own, copied from the fixed publisher source (SIGMA 18 (2022) 093, source label intersection, source lines 995-999); no statement, proof step or formula was rewritten, shortened or removed.\par\medskip
\noindent\textit{Editorial note.} The publisher class sigma is not available here and is replaced by a plain amsart wrapper; the preamble is the authors' own macro block and theorem declarations copied verbatim from source lines 18-101, with no mathematical macro invented, renamed or added.\par\medskip
\noindent\textit{Editorial note.} The excerpt carries no section heading, so the section counter is fixed at 5 and the authors' shared theorem counter at 6 before the body. The counted statement therefore prints as Theorem 5.10, which is its number in the article under the authors' own declaration newtheorem\{Theorem\}\{Theorem\}[section]. The only other environments in the excerpt are the three context blocks that precede it, numbered consecutively below it; their authoritative identifiers are the source labels cohomtwis and poincaredual and the source line ranges recorded with each block.\par\medskip
\noindent\textit{Editorial note.} Context blocks are the source results that the statement and its proof use: the cohomology presentation of the twistor spaces (label cohomtwis), the Poincare dual of the canonical section (label poincaredual), the comparison of the twistor class for two spin-c structures with its proof, and the labelled equation div2. They are marked as context and are not counted as this task's proof.\par\medskip
\noindent\textit{Editorial note.} No figure, table or bibliography entry is referenced by the retained spans, so none is carried; the two displayed formulas of the statement and the display in the proof are the authors' own and are reproduced unchanged.\par\medskip
\noindent\textit{Allowed standard prerequisites (editorial).}\begin{itemize}
\item Projective bundle formula. For a complex vector bundle E over M of rank k, H\textasciicircum{}*(P(E)) is the free H\textasciicircum{}*(M)-algebra on x = c\_1(O\_\{E\textasciicircum{}*\}(1)) subject only to x\textasciicircum{}k + c\_1(E)x\textasciicircum{}\{k-1\} + ... + c\_k(E) = 0.. the authors' own statement of the convention at source lines 168-169. It supplies the two presentations of H\textasciicircum{}*(Z\_+(M)) in the variables x and x' that the retained computation rewrites and multiplies; no projective-bundle input beyond this is assumed.
\item Poincare dual of the canonical section as a Thom class. For the canonical section s of P(E + C) given pointwise by [0:...:0:1], the Poincare dual of s\_*[M] is x\textasciicircum{}n + c\_1x\textasciicircum{}\{n-1\} + ... + c\_n, the image of the Thom class of the normal bundle.. the article's own Lemma poincaredual, source lines 839-846, delivered here as a context span. It converts the two twistor sections into explicit degree-six classes in H\textasciicircum{}*(Z\_+(M)); the retained obligation begins after this identification.
\item Spin-c structures as an H\textasciicircum{}2(M)-torsor. Isomorphism classes of spin-c structures on M form a torsor over H\textasciicircum{}2(M), and if [s'] = a . [s] then the Hopf bundles of the two projective-bundle models satisfy H' = H (x) L with c\_1(L) = a, whence x' = x - a.. the article's own lemma and discussion at source lines 934-940 and 942-972, delivered here as a context span. It is exactly what licenses the substitution x' = x + (c\_1 - c\_1')/2 used in the proof; nothing else about spin-c structures is assumed.
\end{itemize}\medskip
\begin{Corollary}
\label{cohomtwis}
Let $M$ be an oriented Riemannian manifold with an orthogonal almost complex structure $J$ compatible with the orientation. We have the following isomorphisms as algebras over $H^*(M)$:
\begin{enumerate}\itemsep=0pt
\item[$1.$] If $\dim(M)=4$, then
\begin{align*}
&H^*(Z_+(M))\cong H^*(M)[x]/\big(x^2+c_1 x\big), \\
&H^*(Z_{-}(M))\cong H^*(M)[x]/\big(x^2 + c_1 x + c_2\big).
\end{align*}
\item[$2.$] If $\dim(M)=6$, then
\begin{align*}
&H^*(Z_+(M))\cong H^*(M)[x]/\big(x^4 +2c_1x^3 +\big(c_1^2+c_2\big) x^2 + (c_1c_2-c_3)x\big), \\
&H^*(Z_{-}(M))\cong H^*(M)[x]/\big(x^4 + 2c_1x^3+ \big(c_1^2+c_2\big)x^2 + (c_1c_2+c_3)x\big). \end{align*}
\end{enumerate}
\end{Corollary}

\begin{Lemma}\label{poincaredual}
Let $E \xrightarrow{\pi} M$ be a complex vector bundle of rank $n$ over a closed oriented mani\-fold~$M$. Let $s \colon M \to \mathbb{P}(E \oplus \C)$ denote the canonical section given pointwise by $[0: \cdots : 0 : 1]$.
Then the Poincar\'e dual of $s_*[M]$ in $H^*(\mathbb{P}(E\oplus \C)) \cong H^*(M)[x]/\big(x^{n+1} + c_1 x^n + \cdots + c_n x\big)$ is
\[
x^n + c_1x^{n-1} + \cdots + c_n,
\]
where $c_i$ denote the Chern classes of $E$.\footnote{Here the orientation on $\mathbb{P}(E\oplus \C)$ is understood to be the one induced by the orientation on $M$ and the canonical orientation on the fiber $\CP^n$ corresponding to $\langle x^n, \big[\CP^n\big] \rangle = 1$.}
\end{Lemma}

\begin{Lemma}
Let $M$ be an oriented Riemannian six-manifold and let $s,s'\colon M \to Z_+(M)$ be
two almost complex structures on $M$ compatible with the orientation of $M$.
Let $a \in H^2(M)$ be $($the unique cohomology class$)$ such that $[s']=a\cdot [s]$.
Then the canonical isomorphism between the expressions for $H^*(Z_+(M))$ given in Corollary~{\rm \ref{cohomtwis}} in terms of $s$, $s'$ is given by
\[
x' \mapsto x - a.
\]
\end{Lemma}
\begin{proof}
Let $E$ and $E'$ be the positive spinor bundles associated to the ${\rm spin}^{\rm c}$ structures~$[s]$ and~$[s']$, respectively, and let $L$ denote a complex line bundle with $c_1(L)=a$.

The relation $a \cdot [s]=[s']$ implies that there is an isomorphism
\[\psi \colon \ L \otimes E \to E'.\]
The kernel in $T_x M \otimes \C$ of right multiplication by a positive spinor
$v \in E_x$ is the same as the kernel of right multiplication by $\lambda \otimes v \in
L_x \otimes E_x$ for any $\lambda \neq 0$, so the canonical isomorphism given by the
composition
\[
\mathbb{P}(L\otimes E) \xrightarrow{\phi'} Z_+(M) \xrightarrow{\phi^{-1}} \mathbb{P}(E)
\]
is the obvious one sending the line $L \otimes (\C v)$ to $\C v$.

We conclude that on $Z_+(M)$ the tautological (Hopf) line bundles $H$ and $H'$ coming
respectively from the identifications with $\mathbb{P}(E)$ and $\mathbb{P}(E')$ satisfy
the relation
\[
H' \cong H \otimes L.
\]
Since $x'=c_1(H'^*)$ and $x=c_1(H^*)$, this completes the proof.
\end{proof}

Given a ${\rm spin}^{\rm c}$ structure $[s]$ we will write $c_1([s])$ for the Chern class of
the line bundle classified by the projection ${\mathcal B}{\rm Spin}^{\rm c}(2n) \xrightarrow{{\mathcal B}\delta} {\mathcal B}S ^1$.
Note that $c_1( a\cdot [s]) = 2a + c_1([s])$ so the difference class $a \in H^2(M)$
between two almost complex structures $s$, $s'$ is a specific ``square root'' of the
difference $c_1'-c_1$. We can suggestively write
\begin{equation}
\label{div2}
x' = x + \frac{c_1-c_1'}{2},
\end{equation}
which is unambiguous on rational cohomology (or when there is no $2$-torsion in $H^2(M)$).

\begin{Theorem}\label{intersection} Let $J$ and $J'$ be two orthogonal almost complex structures inducing the same orientation on a closed Riemannian six-manifold $M$, corresponding to sections $s$ and $s'$ of the twistor bundle. Denote by $c_i$, $c_i'$ their respective Chern classes. Then the intersection number $s_*[M] \cdot s'_*[M]$ in $Z_+(M)$ is given by \[\int_M \tfrac{1}{8}\big( c_1^3 + c_1^2 c_1' - c_1c_1'^2 - c_1'^3 \big) + \tfrac{1}{2} \big( c_1 c_2' + c_1' c_2' \big) - c_3.\]
\end{Theorem}

\begin{proof}
By Corollary~\ref{cohomtwis}, the presentation of $H^*(Z_+(M))$ with reference to $J$ is given by \[H^*(M)[x]/\big(x^4 + 2c_1x^3 + \big(c_1^2+c_2\big)x^2 + \big(c_1c_2-c_3\big)x\big),\] and with reference to $J'$ it is given by \[H^*(M)[x']/\big((x')^4 + 2c_1'(x')^3 + ((c_1')^2 + c_2')(x')^2 + (c_1'c_2'-c_3')x'\big).\] Combining Corollary~\ref{cohomtwis} with Lemma~\ref{poincaredual}, we have that the Poincar\'e dual of $s_*[M]$ is given by \[x^3 + 2c_1x^2 + \big(c_1^2+c_2\big)x + (c_1c_2-c_3),\] and the Poincar\'e dual of $s'_*[M]$ is given by \[(x')^3 + 2c_1'(x')^2 + \big((c_1')^2 + c_2'\big)x' + c_1'c_2'-c_3'.\] Recall from \eqref{div2} that $x' = x + \tfrac{1}{2}(c_1 - c_1')$. (Our calculation will be insensitive to torsion, so we may in fact work in the rational cohomology ring where there is no ambiguity in writing $\tfrac{1}{2}(c_1 - c_1')$.) Using this to rewrite the expression for $s'_*[M]$, and multiplying with $s_*[M]$, yields \[x^3\big( \tfrac{1}{8}\big( c_1^3 + c_1^2 c_1' - c_1c_1'^2 - c_1'^3 \big) + \tfrac{1}{2} ( c_1 c_2' + c_1' c_2' ) - c_3' \big).\] Pairing with the fundamental class of $Z_+(M)$ gives the result. \end{proof}

\end{document}
